\chapter[1939 --- Domagk]{1939: Gerhard Domagk}
\label{ch:1939_gerharddomagk}

\section*{Main Contribution}

\textit{Discovered that prontosil (a sulfa drug) could cure bacterial infections, launching the era of antibacterial chemotherapy.}\cite{nobel-1939,lecture-1939}

\section{Official Citation}

for the discovery of the antibacterial effect of prontosil

\section{Background and Historical Context}

Gerhard Johannes Paul Domagk was born on October 30, 1895, in Lagow, Brandenburg, Germany (now Łagów, Poland). He grew up in a modest family; his father was a school principal, and his mother was a teacher. Despite limited financial resources, the family placed great value on education, and young Gerhard demonstrated an early aptitude for the natural sciences. He attended the Gymnasium in Berlin, where he excelled in biology and chemistry, and in 1914, he enrolled at the University of Kiel to study medicine.

Domagk's studies were interrupted by World War I, during which he served as a medical officer on the Western Front. The experience of treating soldiers infected with wound infections---particularly gas gangrene and tetanus---left a lasting impression on him and sparked his lifelong interest in infectious disease. After the war, he completed his medical studies at the University of Kiel, earning his doctorate in 1921. He subsequently worked at several German universities, including the University of Greifswald and the University of M\"unster, where he became professor of pathological anatomy in 1925.

The scientific context of the late 1930s was one of urgent medical need. Bacterial infections were among the leading causes of death worldwide, and the arsenal of effective treatments was extremely limited. Sulfur compounds had been used in various medical applications since antiquity, and in the 1930s, several researchers were investigating the antibacterial properties of synthetic dyes. The challenge was to find compounds that could kill bacteria within the body without causing unacceptable toxicity to the patient---a problem that had thwarted many previous attempts at chemotherapy.

It was within this challenging but promising scientific environment that Domagk made his landmark discovery.

\section{The Discovery/Contribution}

Gerhard Domagk received the Nobel Prize in Physiology or Medicine ``for the discovery of the antibacterial effect of prontosil.'' His work demonstrated that a red dye called prontosil had remarkable antibacterial properties in vivo, laying the foundation for the development of sulfonamide antibiotics---the first effective chemotherapeutic agents against bacterial infections.

\subsection{The Discovery of Prontosil's Antibacterial Activity}

Domagk's discovery of the antibacterial properties of prontosil was the result of a systematic research program that combined elements of both chemistry and pathology. Working at the Bayer Research Laboratories in Elberfeld, Germany, Domagk collaborated with chemists Fritz Mietzsch and Joseph Klarer, who synthesized a series of azo dyes---compounds characterized by a nitrogen-nitrogen double bond linking two aromatic rings---and tested them for antibacterial activity.

The key compound, prontosil (4-sulfonamido-2',4'-diaminoazobenzene hydrochloride), was a vivid red dye that Mietzsch and Klarer had synthesized as part of their dye research program. Domagk tested prontosil against a variety of bacterial infections in mice, and the results were notable. Mice infected with lethal doses of streptococci, staphylococci, or pneumococci were completely protected by treatment with prontosil, while untreated control animals died rapidly. The compound was also effective against a wide range of other pathogenic bacteria, including those responsible for wound infections, septicemia, and meningitis.

Domagk published his findings in 1935 in a landmark paper in the \textit{Deutsche Medizinische Wochenschrift}, and the results caused a sensation in the medical world. Here, for the first time, was a synthetic compound that could cure systemic bacterial infections---a feat that had eluded researchers for decades. The discovery of prontosil represented the dawn of the antibiotic era, predating even the clinical use of penicillin by several years.

\subsection{The Mechanism of Action}

One of the most puzzling aspects of prontosil's antibacterial activity was the discrepancy between its in vitro and in vivo effects. When tested directly against bacteria in culture (in vitro), prontosil showed little or no antibacterial activity. Yet when administered to infected animals (in vivo), it was remarkably effective. This paradox was resolved by the work of Jacques Tr\'efou\"el and colleagues in France, who demonstrated in 1935 that prontosil is metabolized in the body to produce sulfanilamide (4-aminobenzenesulfonamide), which is the actual active antibacterial agent. Prontosil is essentially a prodrug---a compound that is inactive in its original form but is converted to an active metabolite within the body.

Sulfanilamide works by a mechanism that was not fully understood until decades later. It competes with para-aminobenzoic acid (PABA) for the active site of the enzyme dihydropteroate synthase, which is essential for the synthesis of folic acid in bacteria. Since bacteria must synthesize their own folic acid (unlike humans, who obtain it from their diet), inhibition of this enzyme is lethal to bacteria but relatively harmless to human cells. This selective toxicity---the ability to harm bacteria while sparing the host---is the fundamental principle of antimicrobial chemotherapy, and Domagk's discovery of prontosil was the first demonstration that this principle could be put into practice.

\subsection{Development and Clinical Application}

Following Domagk's initial publication, the clinical potential of prontosil and its derivatives was rapidly explored. The French team led by Tr\'efou\"el published their work on sulfanilamide in 1935, and within months, the compound was being tested in patients around the world. The results were dramatic: patients with streptococcal infections who had been expected to die recovered rapidly with sulfanilamide treatment. Sulfapyridine, a related compound, proved effective against pneumonia, and sulfadiazine was found to be effective against a wide range of gram-positive and gram-negative bacteria.

The impact of Domagk's discovery on clinical medicine was immediate and significant. Before the introduction of sulfonamides, bacterial pneumonia carried a mortality rate of approximately 30--40\%. With sulfonamide treatment, this rate dropped to below 10\%. Puerperal fever (childbed fever), which had killed countless new mothers, became a treatable condition. Wound infections, which had been a major cause of death in wars throughout history, could now be managed effectively.

\section{Impact on Modern Medicine}

The discovery of prontosil's antibacterial activity inaugurated the era of antibacterial chemotherapy, transforming the practice of medicine in ways that are difficult to overstate. While sulfonamides have been largely supplanted by antibiotics such as penicillin, amoxicillin, and fluoroquinolones for many clinical indications, they continue to play an important role in modern medicine.

\subsection{Continued Clinical Use of Sulfonamides}

Sulfamethoxazole, a derivative of sulfanilamide, remains one of the most widely used antibacterial agents in the world, typically in combination with trimethoprim (co-trimoxazole or Bactrim). This combination is the treatment of choice for Pneumocystis pneumonia, a life-threatening opportunistic infection that is particularly common in patients with HIV/AIDS. The combination is also used to treat urinary tract infections, certain types of diarrhea, and community-acquired pneumonia. Sulfadiazine, another sulfonamide, remains important in the treatment of toxoplasmosis and nocardiosis.

\subsection{The Antimicrobial Resistance Crisis}

Perhaps a major legacy of Domagk's discovery is the ongoing battle against antimicrobial resistance. The widespread use of sulfonamides in the 1930s and 1940s was soon followed by the emergence of sulfonamide-resistant bacteria, providing one of the earliest demonstrations of the evolutionary principle that microbial populations can adapt to selective pressures. This experience was a harbinger of the much larger crisis of antibiotic resistance that now threatens global health.

Understanding the mechanisms of sulfonamide resistance---which include mutations in the dihydropteroate synthase gene, acquisition of alternative enzymes, and efflux of the drug from the bacterial cell---has informed our broader understanding of antimicrobial resistance mechanisms. The genetic studies of sulfonamide resistance were among the earliest demonstrations that resistance genes could be transferred between bacteria on plasmids and other mobile genetic elements, providing a mechanistic explanation for the rapid spread of resistance through bacterial populations.

Today, the World Health Organization considers antimicrobial resistance a notable threats to global health, food security, and development. The principles of antimicrobial stewardship---using antibiotics only when necessary, selecting the appropriate agent and dose, and limiting the duration of therapy---were first articulated in response to the early experience with sulfonamides and remain central to efforts to combat resistance.

\subsection{Prodrug Design}

Domagk's observation that prontosil is a prodrug---converted to its active metabolite within the body---has had far-reaching implications for drug design. The concept of prodrugs is now a central strategy in pharmaceutical development, used to improve the bioavailability, reduce the toxicity, or target the delivery of drugs that would otherwise be unsuitable for clinical use. Many widely used medications are prodrugs, including the antiviral agent valacyclovir (which is converted to acyclovir), the antiplatelet agent clopidogrel (Plavix), and the angiotensin-converting enzyme inhibitor enalapril. The prodrug concept that was first demonstrated with prontosil continues to drive innovation in drug development.

\section{Legacy}

Gerhard Domagk's discovery of the antibacterial properties of prontosil was one of the major medical breakthroughs of the twentieth century. By demonstrating that synthetic chemicals could cure bacterial infections, Domagk opened up an entirely new field of medicine and saved countless lives. The Nobel Prize that was awarded to him in 1939 was, in a sense, a recognition not only of his individual achievement but of the power of systematic scientific research to address the most pressing medical challenges of the day.

It is worth noting that Domagk was forced by the Nazi government to decline the Nobel Prize in 1939, as the German regime considered acceptance of the prize to be an act of defiance. He was not permitted to receive the prize and its accompanying medal until 1947, when he was finally able to travel to Stockholm to accept it in a belated ceremony. This episode serves as a poignant reminder of the ways in which political forces can interfere with the free pursuit of scientific knowledge.

Domagk continued his research at the Bayer Laboratories until his retirement, contributing to the development of several additional antibacterial agents, including thiacetazone, which was used in the treatment of tuberculosis. He died on October 24, 1964, at the age of 68. His legacy endures in the sulfonamide antibiotics that bear the mark of his discovery, in the broader field of antimicrobial chemotherapy that he inaugurated, and in the countless lives that have been saved by the application of his research to clinical medicine.

The story of prontosil and Domagk also illustrates the complex relationship between science and society. The discovery was made by a German scientist working for a German pharmaceutical company during the rise of the Nazi regime, yet its benefits were enjoyed by people around the world regardless of nationality or politics. It is a reminder that scientific knowledge transcends borders and ideologies, and that the pursuit of knowledge for the benefit of humanity is a universal endeavor that deserves support and protection from all quarters.


\section{Modern Developments}

Domagk's discovery of prontosil launched the antibiotic era. Modern antibiotics, including fluoroquinolones, cephalosporins, and carbapenems, have saved countless lives. However, antimicrobial resistance (AMR) is now a global health crisis---the WHO estimates AMR caused 1.27 million deaths in 2019. New approaches to combat AMR include phage therapy, antimicrobial peptides, and CRISPR-based antimicrobials. The COVID-19 pandemic highlighted the importance of maintaining antibiotic pipelines for secondary bacterial infections.


\section{Bibliography}

\begin{thebibliography}{9}
\bibitem{nobel-1939}
Nobel Prize Outreach, ``The Nobel Prize in Physiology or Medicine 1939,''\emph{NobelPrize.org}.\\
\url{https://www.nobelprize.org/prizes/medicine/1939/summary/}

\bibitem{lecture-1939}
Gerhard Domagk, ``Nobel Lecture,''\emph{NobelPrize.org}. Winner-authored primary source; downloadable from the official Nobel Prize archive.\\
\url{https://www.nobelprize.org/prizes/medicine/1939/domagk/lecture/}
\end{thebibliography}
